\documentclass[12pt]{article}
\usepackage[utf8]{inputenc}
\usepackage[T1]{fontenc}
\usepackage{amsmath,amssymb,amsthm}
\usepackage{geometry}
\geometry{margin=1in}
\usepackage{enumitem}
\setlist{nosep}

\linespread{1.1} % Slightly increased line spacing for better readability
% -------------------- Macros --------------------
\newcommand{\N}{\mathbb{N}}
\newcommand{\Z}{\mathbb{Z}}
\newcommand{\Q}{\mathbb{Q}}
\newcommand{\R}{\mathbb{R}}
\newcommand{\C}{\mathbb{C}}
\newcommand{\K}{\mathbb{K}} % used to denote either R or C
\newcommand{\eps}{\varepsilon}
\newcommand{\dd}{\,\mathrm{d}}

\newcommand{\NumberedDefinition}[3]{% #1 number, #2 title, #3 body
\paragraph*{Definition #1 ( #2 ).} #3\par}
\newcommand{\NumberedTheorem}[3]{% #1 number, #2 title, #3 body
\paragraph*{Theorem #1 ( #2 ).} #3\par}
\newcommand{\NumberedProposition}[3]{% #1 number, #2 title, #3 body
\paragraph*{Proposition #1 ( #2 ).} #3\par}
\newcommand{\NumberedLemma}[3]{% #1 number, #2 title, #3 body
\paragraph*{Lemma #1 ( #2 ).} #3\par}
\newcommand{\NumberedCorollary}[3]{% #1 number, #2 title, #3 body
\paragraph*{Corollary #1 ( #2 ).} #3\par}
\newcommand{\NumberedRemark}[3]{% #1 number, #2 title, #3 body
\paragraph*{Remark #1 ( #2 ).} #3\par}
\newcommand{\NumberedExample}[3]{% #1 number, #2 title, #3 body
\paragraph*{Example #1 ( #2 ).} #3\par}

% This chapter translation: Original German "Die Exponentialfunktion" (The Exponential Function)
% Numbering (Definition/Satz) preserved exactly as in source. Do NOT renumber.


% ============================================================================
% 3 Sequences and Series (Folgen und Reihen) -- Original start p.31
% ============================================================================

\begin{document}

\subsection*{3.2. Completeness (Vollständigkeit)}

In the proofs of this chapter we will occasionally use the (complete) induction principle to establish statements for natural numbers.

\NumberedDefinition{3.2.1}{Monotone sequences}{A sequence $(x_n)_{n\in\N}$ in $\R$ is called \emph{monotone increasing} (resp. \emph{monotone decreasing}) if for all $n\in\N$ we have
\[
	x_{n+1} \ge x_n \quad (\text{resp. } x_{n+1} \le x_n ).
\]}

\NumberedTheorem{3.2.2}{Monotone bounded sequences converge}{Every monotone increasing (resp. monotone decreasing) bounded sequence $(x_n)_{n\in\N}$ in $\R$ is convergent, and moreover
\[
	\lim_{n\to\infty} x_n = \sup\{x_n : n\in\N\} \quad\big(\text{resp. }\ \lim_{n\to\infty} x_n = \inf\{x_n : n\in\N\}\big).
\]}

\paragraph{Proof.} We prove the claim for monotone increasing sequences; the decreasing case is analogous.

Assume $(x_n)$ is monotone increasing and bounded. Then the set $\{x_n : n\in\N\}$ is bounded above, hence it has a supremum $s:=\sup\{x_n : n\in\N\}$ by the least upper bound property of $\R$. Let $\eps>0$. By the definition of the supremum there exists $n_\eps\in\N$ with
\[
	s-\eps < x_{n_\eps} \le s.
\]
By monotonicity this inequality holds for all $n\ge n_\eps$, i.e. $|x_n-s|<\eps$ for all $n\ge n_\eps$. Hence $x_n\to s$.
\qed

\bigskip

In the following we compile some statements about natural numbers that can be proved by complete induction and that are partly given as exercise problems. We need the notation: for $k\in\N$ the factorial is defined recursively by
\[
    0! := 1, \qquad k! := k\,(k-1)! \quad (k=1,2,\dots).
\]
For natural numbers $n\ge k$ the binomial coefficient is defined by
\[
    \binom{n}{k} := \frac{n!}{k!\,(n-k)!}.
\]

\paragraph*{Lemma 3.2.3.}
\begin{enumerate}[label=(\arabic*),leftmargin=*]
    \item (Binomial theorem) For all real numbers $x,y$ and all $n\in\N$ we have
    \[
        (x+y)^n = \sum_{k=0}^{n} \binom{n}{k} x^k y^{n-k}.
    \]

    \item For all $k\ge 0$ we have
    \[
        2^{k-1} \le k!.
    \]

    \item (Sums for the geometric series) For every real number $x\ne 1$ and every $n\in\N$ we have
    \[
        \sum_{k=0}^{n} x^k = \frac{1-x^{n+1}}{1-x}.
    \]

    \item (Bernoulli inequality) For all $x>-1$ and $n\ge 0$ we have
    \[
        (1+x)^n \ge 1 + n x.
    \]
\end{enumerate}

\NumberedExample{3.2.4}{Powers and the number $e$}{
\begin{enumerate}[label=(\arabic*),leftmargin=*]
    \item Let $a\in\K$. Then:
    \[
        a^n \xrightarrow[n\to\infty]{} 0 \quad\text{if } |a|<1,\qquad
        a^n \xrightarrow[n\to\infty]{} 1 \quad\text{if } a=1.
    \]
    The sequence $(a^n)_{n\in\N}$ diverges if $|a|\ge 1$ and $a\ne 1$.

    \item Let $a\in\K$ with $|a|<1$ and $r\in\N$. Then
    \[
        n^r a^n \xrightarrow[n\to\infty]{} 0.
    \]

    \item We have
    \[
        \lim_{n\to\infty} \sqrt[n]{n} = 1.
    \]

    \item Let $a\in\R_{>0}$. Then
    \[
        \lim_{n\to\infty} \sqrt[n]{a} = 1.
    \]

    \item The sequence $\bigl((1+\tfrac1n)^n\bigr)_{n\in\N}$ is convergent. Its limit is called \emph{Euler’s number} and is denoted by $e$. We have
    \[
        2<e<3.
    \]

    \item
    \[
        \lim_{n\to\infty} \Bigl(\sum_{k=0}^{n} \frac1{k!}\Bigr) = e.
    \]
\end{enumerate}
}

\paragraph{Proof.}
\begin{enumerate}[label=(\arabic*),leftmargin=*]
    \item Assume that $(a^n)_{n\in\N}$ is convergent. Let
    \[
        \alpha := \lim_{n\to\infty} a^n.
    \]
    Then
    \[
        \alpha = \lim_{n\to\infty} a^n
        = \lim_{n\to\infty} a^{n+1}
        = \lim_{n\to\infty} (a\cdot a^n)
        = \alpha\cdot\alpha.
    \]
    Hence the possibilities are: either $\alpha=0$ and $a$ arbitrary, or $a=1$ and $\alpha$ arbitrary. In the second case it follows that
    \[
        \alpha = \lim_{n\to\infty} a^n = 1.
    \]

    Now let $|a|<1$. Then $(|a|^n)_{n\in\N}$ is monotonically decreasing and bounded, hence convergent by Theorem~3.2.2. Since $a\ne 1$, it follows that
    \[
        \alpha = \lim_{n\to\infty} a^n = 0.
    \]

    Let $|a|=1$ and $a\ne 1$. Suppose $(a^n)_{n\in\N}$ were convergent. Then $\alpha = \lim_{n\to\infty} a^n = 0$, and thus
    \[
        \lim_{n\to\infty} |a|^n = 0,
    \]
    which contradicts $|a|=1$. Hence $(a^n)_{n\in\N}$ diverges.

    Finally, let $|a|>1$. Then, by the first part of the proof,
    \[
        \Bigl(\frac1{|a|}\Bigr)^n \xrightarrow[n\to\infty]{} 0.
    \]
    Thus for every $\eps>0$ there exists $n_\eps\in\N$ such that
    \[
        \frac1{|a|^n} < \eps \qquad\forall n\ge n_\eps,
    \]
    i.e. $|a|^n > 1/\eps$ for all $n\ge n_\eps$. Hence $(a^n)_{n\in\N}$ is unbounded. The claim then follows from Theorem~3.1.8.

    \item Let $a\ne 0$. Set $\alpha := |a|$ and $x_n := n^r \alpha^n$. Then $x_n>0$, and\footnote{The sequence $1+\tfrac1n$ converges to $1$. From Theorem~3.1.19(4) it follows inductively that finite products of convergent sequences converge to the product of the limits. Hence $(1+\tfrac1n)^r$ converges to $1$.}
    \[
        \frac{x_{n+1}}{x_n}
        = \frac{(n+1)^r\alpha^{n+1}}{n^r\alpha^n}
        = \alpha\Bigl(1+\frac1n\Bigr)^r
        \xrightarrow[n\to\infty]{}\alpha.
    \]
    Let $\beta\in(\alpha,1)$. Then there exists $n_0\in\N$ with
    \[
        \frac{x_{n+1}}{x_n}-\alpha
        = \Bigl|\frac{x_{n+1}}{x_n}-\alpha\Bigr|
        < \beta-\alpha \qquad\forall n\ge n_0,
    \]
    and hence $\frac{x_{n+1}}{x_n}<\beta$ for all $n\ge n_0$, i.e.
    \[
        x_{n+1} < \beta x_n \qquad\forall n\ge n_0.
    \]
    By induction it follows that
    \[
        x_n \le x_{n_0}\beta^{\,n-n_0} \qquad\forall n\ge n_0.
    \]
    Thus
    \[
        |n^r a^n| = x_n \le x_{n_0}\beta^{\,n-n_0} \xrightarrow[n\to\infty]{} 0
    \]
    (since $0\le\beta<1$). The assertion follows from Lemma~3.1.18(3).

    \item We show $\sqrt[n]{n}\to 1$. Let $\eps>0$. Because $0<\frac1{1+\eps}<1$, part (2) yields
    \[
        \lim_{n\to\infty} n\Bigl(\frac1{1+\eps}\Bigr)^n = 0
        \qquad\text{(choose $r=1$, $a=\tfrac1{1+\eps}$).}
    \]
    Hence there exists $n_0\in\N$ such that
    \[
        0 < n\Bigl(\frac1{1+\eps}\Bigr)^n < 1 \qquad\forall n\ge n_0 > 1,
    \]
    which implies
    \[
        1 \le n < (1+\eps)^n.
    \]
    From the addendum to this proof it follows that the above estimate implies the inequality
    \[
        1 \le \sqrt[n]{n} < 1+\eps \qquad\forall n\ge n_0.
    \]
    This proves the claim.

    \item Since $a>0$, there exists $n_0\in\N$ with
    \[
        \frac1n < a < n \qquad\forall n\ge n_0.
    \]
    It follows that
    \[
        \frac1{\sqrt[n]{n}} < \sqrt[n]{a} < \sqrt[n]{n} \qquad\forall n\ge n_0.
    \]
    The assertion follows from part (3), Theorem~3.1.19(5), and Theorem~3.1.20.

    \item Let $e_n := \bigl(1+\tfrac1n\bigr)^n$. Show that the sequence $(e_n)_{n\in\N}$ is bounded and monotone. \\
\textbf{Bounded:}

\[
e_n = \Bigl(1 + \frac{1}{n}\Bigr)^n
\overset{\text{Lemma 3.2.3(1)}}{=}
\sum_{k=0}^{n} \binom{n}{k}\frac{1}{n^k}
= 1 + \sum_{k=1}^{n} \binom{n}{k}\frac{1}{n^k}
\]
\[
= 1 + \sum_{k=1}^{n}
\frac{\underbrace{n(n-1)\cdots (n-k+1)}_{k\ \text{factors}}}{k!}\,
\frac{1}{n^k}
\le 1 + \sum_{k=1}^{n} \frac{1}{k!}
\]
\[
{}^{*}\le 1 + \sum_{k=1}^{n} 2^{-(k-1)}
\quad\text{[follows from Lemma 3.2.3(2)]}
\]
\[
= 1 + \sum_{k=0}^{n-1} \Bigl(\frac12\Bigr)^{k}
\overset{\text{Lemma 3.2.3(3)}}{=}
1 + \frac{1 - (1/2)^n}{1 - 1/2}
< 1 + \frac{1}{1/2}
= 3.
\]

\textbf{Monotonicity:}

\[
\frac{e_n}{e_{n-1}}
= \frac{(n+1)^n}{n^n}\cdot\frac{(n-1)^{n-1}}{n^{n-1}}
= \frac{n}{n-1}\,\frac{\bigl((n+1)(n-1)\bigr)^n}{(n^2)^n}
= \frac{n}{n-1}\Bigl(1 - \frac{1}{n^2}\Bigr)^n.
\]
From the Bernoulli inequality (Lemma 3.2.3(4)) it follows that
\[
\frac{e_n}{e_{n-1}}
\ge \frac{n}{n-1}\Bigl(1 - \frac{1}{n}\Bigr)
= 1.
\]
Furthermore, $e_n \ge e_1 = 2$, and thus the assertion follows.

\item Set $x_n := \sum_{k=0}^{n} \frac{1}{k!}$. Monotonicity is obvious. From (*) it follows that $0 \le x_n \le 3$. Hence $(x_n)$ is convergent. Set $e' := \lim_{n\to\infty} x_n$.

Show: $e \le e'$. This follows directly from the estimate (*) above: $e_n \le x_n$ and Theorem~3.1.20.

Show: $e \ge e'$. Let $m\in\N$ be arbitrary but fixed. For $n>m$ we have
\[
    e_n = \Bigl(1 + \frac{1}{n}\Bigr)^n
    = \sum_{k=0}^{n} \binom{n}{k} \frac{1}{n^k}
    > \sum_{k=0}^{m} \binom{n}{k} \frac{1}{n^k}
    =: y_n.
\]
We have $y_n \to x_m$ because
\[
    y_n
    = \sum_{k=0}^{m} \frac{1}{k!} \frac{n(n-1)\cdots (n-k+1)}{\underbrace{n\cdot n\cdots n}_{k\ \text{times}}}
    = \sum_{k=0}^{m} \Bigl(\frac{1}{k!}\Bigr)
    \Bigl(1\cdot\Bigl(1-\frac{1}{n}\Bigr)\Bigl(1-\frac{2}{n}\Bigr)\cdots\Bigl(1-\frac{k-1}{n}\Bigr)\Bigr)
    \xrightarrow[n\to\infty]{}
    \sum_{k=0}^{m} \frac{1}{k!}
    = x_m.
\]

The sequences $(e_n)$ and $(y_n)$ converge, and $e_n > y_n$. Hence $e \ge x_m$ for all $m\in\N$. It follows further that $e \ge e'$.

\[
\Longrightarrow\ e = e'.
\]
\end{enumerate}

\textit{Addendum: (exercise)}
\begin{enumerate}[label=(\arabic*),leftmargin=*]
    \item $0<a<b \Longrightarrow 0<\sqrt[n]{a}<\sqrt[n]{b}$.
    \item $0<a\le b \Longrightarrow 0<\sqrt[n]{a}\le\sqrt[n]{b}$.
\end{enumerate}
\qed

\begin{center}
\textit{Central theorem of this section}
\end{center}

\NumberedTheorem{3.2.5}{Bolzano-Weierstrass}{
Every bounded sequence in $\K$ possesses at least one accumulation point.}

\noindent\textbf{Proof.} Let first $\K = \R$ and $(x_k)_{k \in \N}$ be a bounded sequence in $\R$. Then there exist $A_0 < B_0$ with $A_0 \le x_k \le B_0 \quad \forall k \in \N$.

Starting from $A_0$ and $B_0$, we construct two sequences $(A_n)_{n \in \N}, (B_n)_{n \in \N}$ recursively. Let $A_n, B_n$ already be constructed.

\textbf{Case 1:} There are infinitely many $k \in \N$ with $x_k = \frac{1}{2}(A_n + B_n)$. Then $c := \frac{1}{2}(A_n + B_n)$ is obviously an accumulation point of $(x_k)_{k \in \N}$.

\textbf{Case 2:} Infinitely many sequence terms of $(x_k)_{k \in \N}$ lie in $[A_n, \frac{1}{2}(A_n + B_n)]$. Then we set $A_{n+1} = A_n$ and $B_{n+1} = \frac{1}{2}(A_n + B_n)$.

\textbf{Case 3:} Infinitely many sequence terms of $(x_k)_{k \in \N}$ lie in $]\frac{1}{2}(A_n + B_n), B_n]$. Then we set $A_{n+1} = \frac{1}{2}(A_n + B_n)$ and $B_{n+1} = B_n$.

Obviously: $A_n$ is monotonically increasing and $B_n$ monotonically decreasing and bounded:
\[
A_n, B_n \in [A_0, B_0] \quad \forall n \in \N
\]
Therefore $(A_n)_{n \in \N}, (B_n)_{n \in \N}$ converge. Because $B_n - A_n \le 2^{-n}(B_0 - A_0) \quad \forall n \in \N$ the limits coincide.

Set $c := \lim_{n \to \infty} A_n (= \lim_{n \to \infty} B_n)$. We still show: $c$ is an accumulation point.
Let $\eps > 0$. Then there exists an $n_{\eps_1} \in \N$ and $n_{\eps_2} \in \N$ with
\[
c - \eps < A_n \le c \quad \forall n \ge n_{\eps_1}, \qquad c + \eps > B_n \ge c \quad \forall n \ge n_{\eps_2}
\]
Set $n_\eps = \max\{n_{\eps_1}, n_{\eps_2}\}$. By construction, there are infinitely many $k \in \N$ with
\[
c - \eps < A_{n_\eps} < B_{n_\eps} < c + \eps, \qquad x_k \in [A_{n_\eps}, B_{n_\eps}].
\]
Therefore $c$ is an accumulation point.

Let now $\K = \C$ and $(x_k)_{k \in \N}$ be a bounded sequence in $\C$. Define real sequences:
\[
u_k := \text{Re}(x_k) \quad , \quad v_k := \text{Im}(x_k), \qquad k \in \N.
\]
Because $|\text{Re}(z)| \le |z|, |\text{Im}(z)| \le |z|$, the sequences $(u_k)_{k \in \N}, (v_k)_{k \in \N}$ are bounded. According to Part 1 of the proof, there exists an accumulation point $a \in \R$ of $(u_k)$, i.e., there ex. subsequence $(u_{k_l})_{l \in \N}$ with $u_{k_l} \xrightarrow{l \to \infty} a$.

\textit{Consider:} $(v_{k_l})_{l \in \N}$. $(v_{k_l})$ is bounded, i.e., there ex. accumulation point $b$, i.e., there ex. further subsequence $(v_{k_{l_m}})_{m \in \N}$ with $b = \lim_{m \to \infty} v_{k_{l_m}}$. Because $|z| \le |\text{Re}[z]| + |\text{Im}[z]|$ holds:
\[
\left| x_{k_{l_m}} - (a + ib) \right| \le \underbrace{\left| u_{k_{l_m}} - a \right|}_{\text{null seq.}} + \underbrace{\left| v_{k_{l_m}} - b \right|}_{\text{null seq.}} \xrightarrow{m \to \infty} 0
\]
and $a + ib$ is an accumulation point of $(x_k)_{k \in \N}$. \qed

\NumberedDefinition{3.2.6}{Cauchy sequence}{A sequence $(x_n)_{n\in\N}$ in $\K$ is called a \emph{Cauchy sequence} (abbrev. C.S.) if for every $\eps>0$ there exists $n_0\in\N$ such that
\[
	|x_n - x_m| \le \eps \qquad \forall\, n,m\ge n_0.
\]}

\NumberedTheorem{3.2.7}{Every convergent sequence is Cauchy}{Every convergent sequence in $\K$ is a Cauchy sequence.}

\paragraph{Proof.} Let $\eps>0$ and suppose $x_n\to a$ in $\K$. Then there exists $n_\eps\in\N$ such that $|x_n-a|<\eps/2$ for all $n\ge n_\eps$. For $n,m\ge n_\eps$ we have
\[
	|x_n-x_m| = |(x_n-a)-(x_m-a)| \le |x_n-a| + |x_m-a| < \eps/2+\eps/2 = \eps.
\]
Thus $(x_n)$ is Cauchy. \qed

\NumberedTheorem{3.2.8}{Every Cauchy sequence is bounded}{Every Cauchy sequence in $\K$ is bounded.}

\paragraph{Proof.} By assumption there exists $n_1\in\N$ such that $|x_n-x_m|<1$ for all $n,m\ge n_1$. Define
\[
	C := \max\{ |x_n|+1 : n\le n_1\}.
\]
Then for $n\le n_1$ we have $|x_n|\le C$ by definition, and for $n>n_1$ we estimate
\[
	|x_n| = |x_n-x_{n_1}+x_{n_1}| \le |x_n-x_{n_1}| + |x_{n_1}| < 1 + |x_{n_1}| \le C.
\]
Hence $(x_n)$ is bounded. \qed

\NumberedTheorem{3.2.9}{Cauchy implies convergence in $\K$}{Every Cauchy sequence in $\K$ is convergent.}

\paragraph{Proof.} By Theorem~3.2.8 a Cauchy sequence $(x_n)$ is bounded. Therefore, by Bolzano--Weierstrass (Theorem~3.2.5), $(x_n)$ has an accumulation point $a\in\K$. We claim that $x_n\to a$.

Let $\eps>0$. Since $(x_n)$ is Cauchy, there exists $n_\eps\in\N$ such that $|x_n-x_m|<\eps/2$ for all $n,m\ge n_\eps$. Moreover, because $a$ is an accumulation point, there exists an index $k_\eps\ge n_\eps$ with $|x_{k_\eps}-a|<\eps/2$. Then for all $n\ge n_\eps$ we have
\[
	|x_n-a| \le |x_n-x_{k_\eps}| + |x_{k_\eps}-a| < \eps/2 + \eps/2 = \eps.
\]
Thus $x_n\to a$, proving convergence. \qed

\paragraph{Note.} In the proof of Theorem~3.2.9, the crucial external input was the Bolzano--Weierstrass theorem. The proof of Bolzano--Weierstrass itself used that in $\K$ (via $\R$) bounded monotone sequences converge to the supremum/infimum of their ranges. Equivalently: $\R$ is order complete; i.e., every nonempty bounded subset has a supremum/infimum.

\NumberedExample{3.2.10}{Cauchy in $\Q$ without $\Q$-limit}{There exist Cauchy sequences in $\Q$ which do not converge in $\Q$. For instance, construct rational numbers that approximate $\sqrt{2}$ by bisection on $[1,2]$:
\begin{itemize}[leftmargin=*]
	\item Start with the interval $I_0=[1,2]$. Its midpoint $m_0=3/2$ satisfies $m_0^2=2.25>2$, hence the root of $x^2-2=0$ lies in $[1,3/2]$.
	\item Inductively, having $I_n=[a_n,b_n]$ with $a_n,b_n\in\Q$ and $a_n^2\le 2\le b_n^2$, let $m_n:=\tfrac{a_n+b_n}{2}$. If $m_n^2\le 2$ set $I_{n+1}=[m_n,b_n]$, else set $I_{n+1}=[a_n,m_n]$.
\end{itemize}
Then $a_n,b_n\in\Q$, $I_{n+1}\subset I_n$, and $|I_n|=b_n-a_n=2^{-n}(b_0-a_0)=2^{1-n}$, so $(a_n)$ and $(b_n)$ are Cauchy sequences in $\R$. Any choice of a point $x_n\in I_n\cap\Q$ defines a rational sequence $(x_n)$ with $|x_n-x_m|\le |I_{\min\{n,m\}}|\to 0$; hence $(x_n)$ is Cauchy.

However, in $\Q$ the sequence cannot converge, because its (real) limit must be the unique number $c\in[1,2]$ with $c^2=2$, namely $c=\sqrt{2}\notin\Q$ (cf. Remark~2.3.3). Thus $(x_n)$ is Cauchy in $\Q$ but has no limit in $\Q$; viewed in $\R$ it converges to $\sqrt{2}$.}

\NumberedRemark{3.2.11}{Completion viewpoint}{The real numbers $\R$ can be obtained from $\Q$ by \emph{completion}: one considers the set of Cauchy sequences in $\Q$ and identifies two such sequences if their difference is a null sequence; the resulting equivalence classes form an ordered field that is order complete and isomorphic to $\R$ (Theorem~2.3.4). Under this identification, a rational $q\in\Q$ corresponds to the constant sequence $(q,q,\dots)$, while an irrational such as $\sqrt{2}$ is represented by any rational Cauchy sequence that converges to it in $\R$ (e.g. the bisection sequence above).}

\end{document}